\documentclass[a4paper,10pt]{report}\input{../0.Préambule/Préambule_cours_prof}\begin{document}

\newpage
\section*{Nombres relatifs et enchainement d'opérations}
\addcontentsline{toc}{section}{Nombres relatifs et enchainement d'opérations}

\vspace{-2mm}
\subsection*{Rappel des priorités}
\addcontentsline{toc}{subsection}{Rappel des priorités}

\vspace{-2mm}
\begin{exer1}
Entoure le signe opératoire de l'opération prioritaire. (Il peut y en avoir plusieurs)

\begin{enumerate}[font=\color{exer1}]
\begin{tabularx}{\linewidth}{*{4}{>{\arraybackslash}X}}
\item $O = 252 + 21 \times 41$ &
\item $X = 6,3 - 2,1 \div 7$ &
\item $Y = 3 + 0,3  \times -3$ &
\item $D = 2 \times 2 - 2 \div 2$ \\
\item $A = 17 - 15 \div 3 + 1$ &
\item $B = 50 + 3 + 2 \times 10$ &
\item $L = 0,204 \times 99 - 5,4$ &
\item $E =9 + 12 \times 11 \div 8$ \\
\end{tabularx}
\end{enumerate}
\end{exer1}

\begin{exer2}
Effectue les calculs suivants en soulignant le(s) calcul(s) en cours.

\begin{enumerate}[font=\color{exer2}]
\begin{tabularx}{\linewidth}{*{3}{>{\arraybackslash}X}}
\item $J = 24 + 3 \times 7$ &
\item $O = 720 \div 9 + 4$ &
\item $C = 15 \div 5 - 2$ \\
\item $K = 20 - 0,1 \times 38$ &
\item $E = 60 - 14 + 5 \times 3 + 2$ &
\item $Y = 8 \times 3 - 5 \times 4 \times 0,2$ \\
\end{tabularx}
\end{enumerate}
\end{exer2}

\begin{exer1}
Entoure le signe opératoire de l'opération prioritaire. (Il peut y en avoir plusieurs)

\begin{enumerate}[font=\color{exer1}]
\begin{tabularx}{\linewidth}{*{4}{>{\arraybackslash}X}}
\item $K = (6,2 - 0,1) \div 10$ &
\item $Y = 238 - 4 \times (13 + 27)$ &
\item $S = 5 + (2,8 + 6 \times 1,2)$ &
\item $T = 90 - (2 \times 7 - 7) \times 6$ \\
\item $I = 9 \div 3 + (15 - 6 \div 3)$ &
\item $Q = 34 - ( 104 \div 52 \times 6)$ &
\item $U = (84 - 1) \div (5 + 0,4)$ &
\item $E = 3 \times \left( (1+2) \times 4 - 2 \right)$ \\
\end{tabularx}
\end{enumerate}
\end{exer1}

\begin{exer2}
Effectue les calculs en soulignant le calcul en cours.

\begin{enumerate}[font=\color{exer2}]
\begin{tabularx}{\linewidth}{*{3}{>{\arraybackslash}X}}
\item $A = 25 - (8 - 3) + 1$ &
\item $Z = 25 - 8 - (3 + 1)$ &
\item $Y = 25 - (8 - 3 + 1)$ \\
\item $M = 18 - \left( 4 \times (5 - 3) + 2 \right)$ &
\item $E = 24 \div \left(8 - (3 + 1) \right)$ &
\item $S = \left( 2 + 0,1 \times (5 + 3) \right) \div 4$ \\
\end{tabularx}
\end{enumerate}
\end{exer2}

\begin{exer2}
Calcule les expressions suivantes.

\vspace{-0.75cm}
\begin{enumerate}[font=\color{exer2}]
\begin{tabularx}{\linewidth}{*{3}{>{\arraybackslash}X}}
&&\item $K = 35 - \left( 4 \times (5 + 2) - 7 \right)$ \\
\item $A = 12 \times \left( 32 - (4 + 7) \times 2 \right)$ &
\item $M = (1 + 7) \times \left( 11 - (2 + 3) \right)$ &
\item $I = 12 + \left( (120 - 20) - 2 \times 4 \times 5 \right)$ \\
\end{tabularx}
\end{enumerate}
\end{exer2}

\begin{exer2}
Calcule chacune des expressions suivantes.

\begin{enumerate}[font=\color{exer2}]
\begin{tabularx}{\linewidth}{*{3}{>{\arraybackslash}X}}
\item $A = \dfrac{81}{9} \times 5 - 1$ &
\item $V = \dfrac{45,5}{2} \times 3 - 1$ &
\item $I = \dfrac{27}{2 \times 3} - 1$ \\
\item $S = \dfrac{17 - 5}{3} + 2$ &
\item $E = 7 \times \dfrac{15 \times 4}{3 - 2} + 2 \times 8$ &
\item $Z = \dfrac{13 \times (4 + 7) - 5}{13 - (2 \times 4 + 3}$ \\
\end{tabularx}
\end{enumerate}
\end{exer2}

\vspace{-1mm}
\subsection*{Écrire en une seule expression}
\addcontentsline{toc}{subsection}{Écrire en une seule expression}

\vspace{-2mm}
\begin{exer2}
Voici un programme de calcul 

Applique ce programme à chacun des nombres en écrivant \\une seule expression permettant de trouver le résultat puis calcule :

\begin{enumerate}[font=\color{exer2}]
\begin{tabularx}{\linewidth}{*{4}{>{\arraybackslash}X}}
\item $8$ &
\item $4$ &&\\
\end{tabularx}
\end{enumerate}

\vspace{-38mm}
\begin{flushright}
\arrayrulecolor{black}
\renewcommand{\arraystretch}{0.1}
\begin{tabular}{|p{5cm}|}
\hline
\begin{verbatim} Choisis un nombre \end{verbatim}
\begin{verbatim} Ajoute -8 \end{verbatim}
\begin{verbatim} Multiplie par 3 \end{verbatim}
\begin{verbatim} Divise par -2 \end{verbatim}
\\\hline
\end{tabular}
\end{flushright}
\end{exer2}

\vspace{-2mm}
\begin{exer2}
Maya achète $5$ pots de confitures à $1,80$\euro pièce et \\$12$ baguettes de pain à $0,70$\euro pièce. 

\noindent Quel est le prix total qu'elle doit payer ?
\end{exer2}

\begin{exer3}
Voici un programme de calcul 

Applique ce programme à chacun des nombres en écrivant \\une seule expression permettant de trouver le résultat puis calcule :

\begin{enumerate}[font=\color{exer3}]
\begin{tabularx}{\linewidth}{*{4}{>{\arraybackslash}X}}
\item $5$ &
\item $-5$ &&\\
\end{tabularx}
\end{enumerate}

\vspace{-38mm}
\begin{flushright}
\arrayrulecolor{black}
\renewcommand{\arraystretch}{0.1}
\begin{tabular}{|p{6cm}|}
\hline
\begin{verbatim} Choisis un nombre \end{verbatim}
\begin{verbatim} Calcule son carre \end{verbatim}
\begin{verbatim} Multiplie par (-10) \end{verbatim}
\begin{verbatim} Soustraire le nombre de depart \end{verbatim}
\\\hline
\end{tabular}
\end{flushright}
\end{exer3}

\begin{exer3}
Une famille de $4$ personnes produit en moyenne $1,42$t de déchets par an. La population française est estimé à environ $67,1$ millions de personnes. 

\noindent Calcul, en kg, la masse de déchets produite de France chaque année. On rappelle que $1$t = $\np{1000}$kg.
\end{exer3}


\end{document}